\chapter{Problem Statement}
\label{chap:problem_statement}

\section{The object of study}

We study a \emph{greedy, zero-shot, tactic-by-tactic} theorem proving agent.
At each step the agent receives the current Lean~4 proof state and the full history
of tactics it has tried; it proposes exactly one tactic; Lean either closes the proof,
advances the goal, leaves it unchanged, or returns an error.
The agent has no lookahead, no backtracking, and no access to the Mathlib source
beyond what the language model learned during pretraining.

Formally, the agent defines a policy
\[
  \pi : (\text{goal state}, \text{tactic history}) \;\longrightarrow\; \text{tactic string}
\]
instantiated by a single frozen language model --- NVIDIA Nemotron-3-Ultra-550B
or Groq \texttt{gpt-oss-20b} --- queried via a standard chat completion API.
The policy is evaluated on 837 theorems drawn from real Mathlib,
stratified into three difficulty tiers (Easy: 310, Medium: 378, Hard: 149).

\section{Research questions}

Five questions structure the study.

\paragraph{RQ1 --- Capability ceiling.}
What pass rate does the agent achieve on Mathlib Easy, Medium, and Hard theorems?
Are the tier labels calibrated to agent difficulty, or do they reflect human difficulty alone?

We report two variants, and the distinction is not cosmetic.
The \emph{headline} pass rate is proved over all theorems attempted.
The \emph{reachable-only} pass rate is proved over those the model actually got to
attempt --- that is, theorems the agent either proved or failed to prove within its
step budget.
Excluded from the reachable denominator are theorems lost before any mathematics
happened: those whose statement failed to elaborate in Lean
(\texttt{FAILED\_ELABORATION}, a dataset-quality signal), and those where the model
was never reached at all because the provider was down, rate-limiting, or returning
errors (\texttt{FAILED\_LLM\_ERROR}, \texttt{FAILED\_CONNECTIVITY\_OUTAGE}).

The second exclusion is the load-bearing one.
Across the two Easy-50 runs conducted so far, 28 of 50 and 22 of 50 theorems
respectively were lost to provider failure or elaboration --- so the headline figure
divides successes by a denominator dominated by infrastructure the study does not
control, and measures API uptime as much as model capability.
Reporting both is what keeps that visible rather than buried:

\begin{center}
\begin{tabular}{@{}lcc@{}}
\hline
\textbf{Run (Easy, $n=50$)} & \textbf{Headline} & \textbf{Reachable-only} \\
\hline
Nemotron-3-Ultra-550B & 7/50 = 14.0\% & 7/22 = 31.8\% \\
Groq \texttt{gpt-oss-20b} & 8/50 = 16.0\% & 8/28 = 28.6\% \\
\hline
\end{tabular}
\end{center}

Wilson score intervals accompany every rate; at these sample sizes they are wide, and
the two models are not distinguishable from one another.

\paragraph{RQ2 --- Failure mode taxonomy.}
When the agent fails, what is the proximate cause?
We define five mutually exclusive failure modes, detectable from the JSONL event log
and assigned first-match-wins in the order below:
\begin{itemize}
  \item \textsc{No-Progress Loop} --- the final three steps all report the
        \texttt{no\_progress} outcome: Lean accepted the tactic but the proof state
        did not change.
  \item \textsc{Error Loop} --- the same Lean message is returned by the last three
        steps that produced any message at all.
  \item \textsc{Hallucinated Lemma} --- Lean reports ``unknown identifier'' or
        ``unknown constant'' while the tactic names a capitalised symbol. Known
        library roots (\texttt{Nat}, \texttt{Int}, \texttt{Real}, \texttt{List},
        \texttt{Set}, \texttt{Finset}, \texttt{Bool}, \texttt{Prop}, \texttt{Type},
        \texttt{Sort}) are whitelisted, so \texttt{exact Nat.add\_comm \_ \_} is not
        misread as an invention.
  \item \textsc{Syntax Error Loop} --- three or more steps return messages mentioning
        ``expected'' or ``parse error''.
  \item \textsc{Budget Exhaustion (Mixed)} --- none of the above;
        the agent made some progress but ran out of steps.
\end{itemize}
The relative frequency of these modes determines which improvement ---
better prompting, search, tool use, or fine-tuning --- would have the largest effect.

Failures caused by the provider rather than the model are excluded before
classification. Without that filter a run of identical \texttt{APIConnectionError}
messages classifies as \textsc{Error Loop}, which would attribute an outage to the
agent's reasoning.

\paragraph{RQ3 --- Step economy.}
How many steps does a successful proof require?
Is there a step threshold beyond which the probability of success effectively vanishes?
A survival curve --- $P(\text{proved} \mid \text{still alive after } k \text{ steps})$ ---
as a function of step number motivates or rules out early stopping heuristics.

\paragraph{RQ4 --- History utilisation.}
Does the agent read and act on its tactic history?
We measure three proxies: the fraction of steps where the model repeats a tactic it already tried
(repetition rate), whether the \emph{advanced} outcome rate increases with step number
(which would indicate the model is learning within the episode),
and how quickly the goal state shrinks (measured by proof-state character count).
The history shown to the model is itself logged per step, so these measures are taken
against what the model actually saw rather than what it was assumed to have seen.

\paragraph{RQ5 --- Tactic distribution.}
Which tactics does the model use, and which ones work?
We separate \emph{usage frequency} from \emph{success rate} (fraction of uses that occur
as the final step of a proved proof) and cluster tactics into reliable closers,
progress makers, and dead ends.

\section{What we do not study}

This study does not optimise pass rate.
We do not implement tree search, self-consistency sampling, fine-tuning,
or any auxiliary tool beyond the Lean REPL.
We do not compare against competition-mathematics benchmarks such as miniF2F or PutnamBench.
The agent is deliberately minimal: its failure modes are therefore maximally interpretable,
not masked by engineering that fixes them before they can be observed.

\section{Success criteria}

The study succeeds if it produces:
\begin{enumerate}
  \item Pass rates with confidence intervals for at least the Easy tier,
        with the headline/reachable distinction clearly stated.
  \item A classified failure breakdown for all failed proofs,
        with at least one documented example per mode.
  \item Enough step-level data to plot the survival curve and
        the repetition-rate statistic.
  \item A tactic frequency and precision table covering all tactics observed.
\end{enumerate}
Medium and Hard results would strengthen RQ1 but are not required for the
other four research questions, which are answerable from the Easy runs alone.
